\paragraph{Conditional vertex-density theorem.}
In the finite-$M$ Moran model, consider the balanced quartet with neutral
dynamics, zero background switching, equal finite copy counts, a fixed
interior initial count, and zero lower-branch durations. For
$z\in\{0011,0101,0110\}$, let $j(z)$ denote the paired topology selected by
$z$. Then there is a finite-state transient survival function $S_z(T)$ such
that
\[
P(Q_i\mid Z=z)=S_z(T)/3+(1-S_z(T))\mathbf 1\{i=j(z)\}.
\]
Moreover $S_z(T)\to0$ as $T\to\infty$, so each of the three simplex vertices
is in the closure of the union of arrangement-conditional Moran quartet laws.

\begin{proof}[Proof sketch]
With zero switching, the paired pattern fixes which two background-homogeneous
pairs can coalesce. Any such coalescence resolves the paired topology. The
unresolved states form a finite killed Moran/genealogy chain; for an interior
initial count every unresolved state has an eventual exit, so its survival
probability tends to zero. At the older boundary, the residual quartet law is
uniform. Conditioning on the positive terminal-pattern mass gives the stated
formula.
\end{proof}

The theorem does not assert surjectivity of the interior. Exact finite-M
coverage for a fixed pattern and for the union over patterns remains open.
